\documentclass[11pt]{article}
   
\usepackage[utf8]{inputenc}
\usepackage[spanish]{babel}
\usepackage[T1]{fontenc}
\usepackage{amsmath, amssymb}
\usepackage{geometry}
\geometry{a4paper, margin=2.5cm}


\title{Relación de Ejercicios: Resolución numérica de ecuaciones no lineales} 
\author{Grado en Ingeniería Informática y Matemáticas \\ Métodos Numéricos II} 
\date{Curso 25/26} 

\begin{document}

\maketitle

\section*{Problema 16: Método de Steffensen} 

El presente problema pretende estudiar el método de Steffensen para la resolución de ecuaciones no lineales. % 

Sea $f:[a,b] \longrightarrow \mathbb{R}$ una función real suficientemente regular y sea $s \in (a,b)$ la única raíz de la ecuación $f(x)=0$. Supongamos además que $f'(x)$ no se anula en $[a,b]$. % [cite: 91]
Se define la función $g:[a,b] \longrightarrow \mathbb{R}$ mediante la expresión: 

\[
g(x) = \begin{cases} 
      \dfrac{f(x+f(x))-f(x)}{f(x)} & x \neq s \\ 
      \rule{0pt}{4ex} \displaystyle \lim_{x \to s} \frac{f(x+f(x))-f(x)}{f(x)} & x = s 
   \end{cases}
\] 

\begin{enumerate}
    \item[a)] Utilice la regla de L'Hôpital para demostrar que $g(s)=f'(s)$. 
    
    \item[b)] Demuestre que $g(x) \neq 0$ para todo $x \in [a,b]$. 
    
    \item[c)] Se define la función $F:[a,b] \longrightarrow \mathbb{R}$ mediante: 
    \[
    F(x) = x - \frac{f(x)}{g(x)}
    \] 
    Demuestre que $F(s) = F'(s) = 0$. 
    
    \item[d)] A partir de los apartados anteriores, encuentre condiciones suficientes para que el método de iteración funcional definido por: % [cite: 103]
    \[
    x_{n+1} = x_{n} - f(x_{n}) \cdot \frac{f(x_{n})}{f(x_{n}+f(x_{n})) - f(x_{n})}
    \] 
    converja localmente a $s$, con velocidad de convergencia al menos cuadrática. 
\end{enumerate}

\begin{enumerate}
    \item [a)]
    Para calcular $g(s)$ debemos evaluar el límite:
\[
g(s) = \lim_{x \to s} \frac{f(x+f(x))-f(x)}{f(x)} = \frac{0}{0}
\]
    
\end{enumerate}

\end{document}